\documentclass[11pt]{article}

\usepackage[T1]{fontenc}
\usepackage[utf8]{inputenc}
\usepackage{lmodern}
\usepackage[margin=1in]{geometry}
\usepackage{amsmath,amssymb,amsthm,mathtools}
\usepackage{hyperref}
\usepackage{microtype}

\hypersetup{colorlinks=true,linkcolor=blue,urlcolor=blue,citecolor=blue}

\newtheorem{theorem}{Theorem}
\newtheorem{lemma}{Lemma}
\newtheorem{remark}{Remark}
\newcommand{\dd}{\mathrm{d}}
\newcommand{\R}{\mathbb{R}}
\newcommand{\cH}{\mathcal{H}}
\newcommand{\cL}{\mathcal{L}}
\newcommand{\bX}{\mathbf{X}}
\newcommand{\bV}{\mathbf{V}}
\newcommand{\ba}{\mathbf{a}}
\newcommand{\bb}{\mathbf{b}}
\newcommand{\bl}{\boldsymbol{\lambda}}
\newcommand{\bp}{\mathbf{p}}
\newcommand{\cQ}{\mathcal{Q}}

\title{A Complete Constraint Count for the Regular Higher-Derivative S004 Complex Scalar}
\author{Richard J. Reyes}
\date{October 26, 2026}

\begin{document}
\maketitle

\begin{abstract}
For the specified nonzero-field, fixed-flat-background sector of the corrected
S004 complex-scalar action, the Hessian with respect to the two real
accelerations is $-2\kappa F I_2$, with nonzero determinant
$4\kappa^2F^2$. An explicit first-order Dirac reduction produces four
second-class primary constraints and no secondary or first-class scalar
constraints. Starting from twelve first-order canonical variables, the
reduced phase space therefore has eight real dimensions per spatial point,
corresponding to four real configuration degrees of freedom. Two degrees
are additional to those of an ordinary complex scalar. The conserved
global $U(1)$ charge is not a local gauge constraint. The result applies
only to the specified regular fixed-background sector; it does not
analyze an unspecified metric or gauge completion or the degenerate vacuum.
\end{abstract}

\section{Hypotheses and action}

Use the rescaled time coordinate $\tau=c_0t$ and the flat metric of
signature $(-,+,\ldots,+)$. Write
\begin{equation}
\psi=X_1+iX_2,\qquad \bX=(X_1,X_2),\qquad
\rho=|\psi|^2=|\bX|^2.
\end{equation}
Define
\begin{equation}
D(\rho)=\rho+\varepsilon^2e^{-2\alpha\rho},\qquad
F(\rho)=\frac{\rho}{D(\rho)^2}.
\end{equation}
Assume
\begin{equation}\label{eq:hyp}
\varepsilon>0,\qquad \kappa>0,\qquad\theta>0,
\qquad\gamma^2<4\kappa\theta,
\end{equation}
and restrict the canonical analysis to a smooth field-space patch
$\rho>0$ (bounded away from zero when uniform regularity is required).
Then $F(\rho)>0$. The scalar background metric and foliation are fixed;
no independent metric, connection, or gauge field is included.

The original S004 scalar density is
\begin{equation}\label{eq:original}
\cL_{\mathrm{orig}}=
\partial_\mu\bar\psi\,\partial^\mu\psi-V(\rho)
+F(\rho)\left[
\kappa|\Box\psi|^2+\theta|\Delta\psi|^2
+\gamma\operatorname{Re}(\Box\psi\,\overline{\Delta\psi})
\right].
\end{equation}
We use $\cL_{\mathrm{phys}}=-\cL_{\mathrm{orig}}$, which
orients the ordinary two-derivative kinetic term positively. An overall
sign reversal does not alter the constraint count.

Put
\begin{equation}
\bV=\partial_\tau\bX,\quad
\ba=\partial_\tau^2\bX,\quad
\bb=\Delta\bX,\quad
A=2\kappa+\gamma,\quad C=\kappa+\theta+\gamma.
\end{equation}
Because $\Box\psi=-\ba+\bb$ in real components,
\begin{align}
|\Box\psi|^2&=|\ba|^2-2\ba\cdot\bb+|\bb|^2,\\
\operatorname{Re}(\Box\psi\,\overline{\Delta\psi})
&=-\ba\cdot\bb+|\bb|^2.
\end{align}
It follows that
\begin{equation}\label{eq:realL}
\boxed{\cL_{\mathrm{phys}}=|\bV|^2-|\nabla\bX|^2+V(\rho)
-F(\rho)\left(\kappa|\ba|^2-A\ba\cdot\bb+C|\bb|^2\right).}
\end{equation}
We assume boundary conditions sufficient to justify the integrations by
parts and the Hamiltonian manipulations below (e.g., periodic spatial data
or suitably decaying variations with vanishing boundary flux).

\section{The acceleration Legendre map}

\begin{lemma}[Regular acceleration Hessian]\label{lem:hessian}
On the patch $\rho>0$, the real two-by-two acceleration Hessian obeys
\begin{equation}\label{eq:hessian}
H_{\mathrm{acc}}=
\left(\frac{\partial^2\cL_{\mathrm{phys}}}
{\partial a_i\partial a_j}\right)_{i,j=1}^{2}
=-2\kappa F(\rho)I_2,
\qquad \det H_{\mathrm{acc}}=4\kappa^2F(\rho)^2>0.
\end{equation}
Thus the higher-time Legendre map is pointwise invertible.
\end{lemma}
\begin{proof}
Differentiating equation~\eqref{eq:realL} with $\bX$, $\bV$,
and $\bb$ fixed gives
\begin{equation}
\frac{\partial\cL_{\mathrm{phys}}}{\partial a_i}
=-F(2\kappa a_i-Ab_i).
\end{equation}
Since $F$ depends on $\bX$, not on $\ba$, one more derivative gives
$\partial^2\cL_{\mathrm{phys}}/\partial a_i\partial a_j
=-2\kappa F\delta_{ij}$. The determinant is positive under
\eqref{eq:hyp} and $\rho>0$.
\end{proof}

\section{First-order Dirac reduction}

Introduce $\bV$ as an independent configuration variable and a
real two-component multiplier $\bl$ imposing $\dot\bX=\bV$:
\begin{equation}\label{eq:auxL}
\cL_1=\cL_{\mathrm{phys}}(\bX,\bV,\dot\bV,\Delta\bX)
+\bl\cdot(\dot\bX-\bV).
\end{equation}
The six real first-order configuration coordinates are
$(\bX,\bV,\bl)$, with conjugate momenta $(\bp_X,\bp_V,\bp_\lambda)$.
The momentum equations are
\begin{equation}\label{eq:momenta}
\bp_X=\bl,\qquad
\bp_\lambda=0,\qquad
\bp_V=-F(2\kappa\dot\bV-A\bb).
\end{equation}
By Lemma~\ref{lem:hessian}, the final relation is solvable:
\begin{equation}\label{eq:invert}
\dot\bV=\frac{A}{2\kappa}\bb-\frac{\bp_V}{2\kappa F}.
\end{equation}
Hence no primary constraint arises from $\bp_V$. The multiplier sector
provides exactly four primary constraints,
\begin{equation}\label{eq:constraints}
\Phi_i=p_{X_i}-\lambda_i\approx0,\qquad
\Xi_i=p_{\lambda_i}\approx0,\qquad i=1,2.
\end{equation}

\begin{lemma}[All four primary constraints are second class]
The Dirac matrix of the four constraints in \eqref{eq:constraints}
has rank four. In particular, no nonzero first-class linear combination
of these primary constraints exists.
\end{lemma}
\begin{proof}
The equal-time canonical brackets are
\begin{equation}
\{X_i(\mathbf y),p_{X_j}(\mathbf z)\}
=\{V_i(\mathbf y),p_{V_j}(\mathbf z)\}
=\{\lambda_i(\mathbf y),p_{\lambda_j}(\mathbf z)\}
=\delta_{ij}\delta(\mathbf y-\mathbf z).
\end{equation}
The nonzero brackets among the constraints are
\begin{equation}
\{\Phi_i(\mathbf y),\Xi_j(\mathbf z)\}
=-\delta_{ij}\delta(\mathbf y-\mathbf z).
\end{equation}
Suppressing the common spatial delta distribution, the matrix in the
order $(\Phi_1,\Phi_2,\Xi_1,\Xi_2)$ is
\begin{equation}\label{eq:diracmatrix}
\mathcal C=\begin{pmatrix}
0&-I_2\\ I_2&0
\end{pmatrix},\qquad\det\mathcal C=1.
\end{equation}
Thus all four constraints are second class.
\end{proof}

\section{Preservation and termination of the Dirac algorithm}

The partial Legendre transform in the acceleration variable, using
\eqref{eq:invert}, gives
\begin{equation}\label{eq:H0}
\boxed{
\cH_0=-|\bV|^2+|\nabla\bX|^2-V(\rho)
-\frac{|\bp_V|^2}{4\kappa F}
+\frac{A}{2\kappa}\bp_V\cdot\bb
+B_\perp F|\bb|^2,}
\end{equation}
where
\begin{equation}
B_\perp=C-\frac{A^2}{4\kappa}
=\theta-\frac{\gamma^2}{4\kappa}>0.
\end{equation}
The total Hamiltonian takes the form
\begin{equation}\label{eq:HT}
H_T=\int\left[\cH_0+\bl\cdot\bV
+\mathbf u\cdot\boldsymbol\Phi
+\mathbf w\cdot\boldsymbol\Xi\right]\dd^d\mathbf y.
\end{equation}
For each $i=1,2$, preservation of $\Xi_i$ yields
\begin{equation}
\dot\Xi_i=\{\Xi_i,H_T\}=u_i-V_i\approx0,
\qquad\Longrightarrow\qquad u_i=V_i.
\end{equation}
Preservation of $\Phi_i$ yields
\begin{equation}
\dot\Phi_i=\{\Phi_i,H_T\}
=-\frac{\delta H_0}{\delta X_i}-w_i\approx0,
\qquad\Longrightarrow\qquad
w_i=-\frac{\delta H_0}{\delta X_i},
\end{equation}
with $H_0=\int\cH_0\dd^d\mathbf y$. The functional derivative
includes all spatial-derivative terms.
Both preservation conditions fix Lagrange multipliers; neither imposes
a new condition on $(\bX,\bV,\bp_X,\bp_V)$. There are therefore
no secondary constraints, and the Dirac algorithm terminates.

\section{Degree-of-freedom theorem}

\begin{theorem}[Four regular real configuration degrees of freedom]
\label{thm:main}
Under the hypotheses in Section~1, the fixed-background corrected S004
complex-scalar sector has four local real configuration degrees of freedom
per spatial point. Two are additional higher-time degrees relative to an
ordinary complex scalar. There are four second-class scalar constraints,
no first-class scalar constraints, and no secondary scalar constraints.
\end{theorem}
\begin{proof}
The auxiliary first-order formulation contains six real configuration
coordinates, hence twelve real canonical phase-space dimensions per spatial
point. The invertible acceleration Hessian introduces no extra primary
constraint; equations~\eqref{eq:constraints}--\eqref{eq:diracmatrix}
give four second-class constraints. The preservation equations above show
that no secondary constraints arise and that there are no first-class
scalar constraints. Dirac's counting formula therefore gives
\begin{equation}
N_{\mathrm{conf}}=
\frac{N_{\mathrm{phase}}-2N_{\mathrm{first}}-N_{\mathrm{second}}}{2}
=\frac{12-2(0)-4}{2}=4.
\end{equation}
Equivalently, eliminating $(\bl,\bp_\lambda)$ leaves four
unconstrained canonical pairs,
\begin{equation}
(X_1,p_{X_1}),\quad (X_2,p_{X_2}),\quad
(V_1,p_{V_1}),\quad (V_2,p_{V_2}),
\end{equation}
so the reduced phase space has eight real dimensions per spatial point.
The two $\bV$ components constitute two additional local higher-time
configuration degrees of freedom.
\end{proof}

\section{Canonical energy and global phase charge}

On the reduced phase space, the exact canonical Hamiltonian is
\begin{equation}\label{eq:Hphys}
\boxed{\begin{aligned}
H_{\mathrm{phys}}=\int\Bigg[
&\bp_X\cdot\bV-|\bV|^2+|\nabla\bX|^2-V(\rho)\\
&-\frac{|\bp_V|^2}{4\kappa F}
+\frac{A}{2\kappa}\bp_V\cdot\Delta\bX
+B_\perp F|\Delta\bX|^2
\Bigg]\dd^d\mathbf y.
\end{aligned}}
\end{equation}
The term $\bp_X\cdot\bV$ is linear in an unconstrained momentum;
moreover, $-|\bp_V|^2/(4\kappa F)$ is negative for nonzero $\bp_V$.
These are explicit manifestations of the unbounded-energy obstruction
in this nondegenerate higher-time sector. This energy observation is
separate from, and not required for, the counting theorem.

Let $J=\begin{psmallmatrix}0&-1\\1&0\end{psmallmatrix}$.
The constant phase transformation $\psi\mapsto e^{i\eta}\psi$
acts as $\delta\bX=\eta J\bX$ and $\delta\bV=\eta J\bV$.
Its canonical generator is the conserved global charge
\begin{equation}\label{eq:charge}
\boxed{\cQ=\int\left(\bp_X\cdot J\bX
+\bp_V\cdot J\bV\right)\dd^d\mathbf y.}
\end{equation}
Conservation presupposes sufficient solution regularity, integrability,
and vanishing boundary flux. Importantly, a \emph{constant} $U(1)$
parameter is not a free spacetime-dependent gauge parameter. The charge
\eqref{eq:charge} is not a local first-class constraint, and fixing its
value does not eliminate the two extra local degrees of freedom.

\section{Limitations}

\begin{remark}[Degenerate vacuum]
At $\psi=0$ one has $F(0)=0$, so the acceleration Hessian loses rank.
The regular-sector Legendre map is then inapplicable. A separate singular
constraint analysis is required at that point; it cannot be used to
eliminate modes throughout neighboring regular nonzero-field patches.
\end{remark}

\begin{remark}[Different theories and completions]
The theorem concerns the literal corrected S004 scalar sector with
fixed flat metric and fixed foliation. It does not supply or count
constraints of an unspecified dynamical-metric, connection, or gauge
completion. The $\kappa=0$ case likewise requires a separate analysis.
For global infinite-domain solutions, one must additionally control
the behavior of $F$ and appropriate canonical function spaces; the
count above is a local regular-sector count.
\end{remark}

\section*{Conclusion}

For $\kappa>0$ and $F(\rho)>0$,
\begin{equation}
\boxed{H_{\mathrm{acc}}=-2\kappa F I_2,\quad
\det H_{\mathrm{acc}}=4\kappa^2F^2>0,\quad
N_{\mathrm{phase,red}}=12-4=8,\quad
N_{\mathrm{conf}}=4.}
\end{equation}
The two extra real higher-time configuration degrees of freedom survive
all scalar constraints of the specified fixed-background system. A
global $U(1)$ charge does not remove them.

\begin{thebibliography}{9}
\bibitem{S004audit}
R.~J. Reyes, ``S004 physical-health audit: constraints, ghost and
ultraviolet obstruction,'' October 20, 2026,
\url{https://github.com/rickyjreyes/geometry_of_resonance/blob/main/docs/foundational/GHOST_AND_CONSTRAINT_ANALYSIS.md}.

\bibitem{Woodard}
R.~P. Woodard, ``The Theorem of Ostrogradsky,''
\emph{Scholarpedia} \textbf{10}(8), 32243 (2015),
\url{https://arxiv.org/abs/1506.02210}.
\end{thebibliography}

\end{document}
